Considera la funció
\[
f(x)=x^2-4x+1 .
\]

\begin{apartats}

\apartat[1,25]{0,75}
Calcula la taxa de variació mitjana de $f$ als intervals $[1,3]$ i $[3,5]$, i interpreta'n el
resultat.

\begin{solucio}
$f(1)=-2$, $f(3)=-2$ i $f(5)=6$.\\
$\mathrm{TVM}[1,3]=\dfrac{f(3)-f(1)}{3-1}=\dfrac{-2-(-2)}{2}=0$: entre $x=1$ i $x=3$ la
funció acaba on havia començat, i en mitjana no creix ni decreix.\\
$\mathrm{TVM}[3,5]=\dfrac{f(5)-f(3)}{5-3}=\dfrac{6-(-2)}{2}=4$: en aquest interval la funció
creix, en mitjana, $4$ unitats per cada unitat de $x$.
\end{solucio}

\apartat[1,25]{1}
Calcula $f'(3)$ aplicant la definició de derivada.

\begin{solucio}
\[
f'(3)=\lim_{h\to0}\frac{f(3+h)-f(3)}{h}
=\lim_{h\to0}\frac{(3+h)^2-4(3+h)+1-(-2)}{h}
=\lim_{h\to0}\frac{h^2+2h}{h}=\lim_{h\to0}(h+2)=2 .
\]
\end{solucio}

\begin{nomesllarg}
\apartat{0,75}
Calcula la taxa de variació mitjana de $f$ a l'interval $[3;\,3{,}1]$ i compara-la amb
$f'(3)$. Com s'explica el resultat?

\begin{solucio}
$f(3{,}1)=9{,}61-12{,}4+1=-1{,}79$, i
$\mathrm{TVM}[3;\,3{,}1]=\dfrac{-1{,}79-(-2)}{0{,}1}=\dfrac{0{,}21}{0{,}1}=2{,}1$.\\
És molt a prop de $f'(3)=2$: la derivada és el límit de les taxes de variació mitjanes als
intervals $[3,3+h]$ quan $h\to0$. Aquí, $\mathrm{TVM}[3,3+h]=h+2$, que per a $h=0{,}1$ val
$2{,}1$.
\end{solucio}
\end{nomesllarg}

\end{apartats}
